
\Needspace{12\baselineskip}
\subsubsection{Opérateurs des \(13\) multisections}

Les nombres de familles indiqués dans chaque fiche sont des cardinalités locales du
support de cette multisection. Ils ne s'additionnent pas entre fiches : le quotient des
56 familles de parcours et le quotient des 13 multisections sont deux partitions
distinctes de l'atlas de 81 cellules, reliées par les applications de projection, et non
par une décomposition disjointe commune.

\begingroup
\small
\setlength{\parskip}{.12\baselineskip}
\setlist[itemize]{leftmargin=1.65em,itemsep=0pt,topsep=.1em,
  parsep=0pt,partopsep=0pt}
\let\texttt\textnormal
Les fiches suivantes présentent l’opérateur dans son codomaine natif. Chaque spectre est reproduit intégralement ; aucune liste de valeurs propres n’est remplacée par une coordonnée de masse.
\Needspace{10\baselineskip}
\noindent\textbf{secteur leptonique chargé, profondeur G0 (centre).} \(\mathcal B_{\ell,\mathrm{G0}}\).\par
\noindent\textbf{Domaine de fibre.} Rang 1; cellule 41; sans orbite chromatique; 1 famille de parcours.
\noindent\textbf{Spectre complet de \(K_D\).}
\begin{itemize}
\item \texttt{\detokenize{(80,54,6)x1}}
\end{itemize}
\noindent\textbf{Spectre complet de \(K_\Omega\).}
\begin{itemize}
\item \texttt{\detokenize{(80,54,6)x1}}
\end{itemize}
\noindent\textbf{Conclusion dans la fibre.} Opérateur diagonal exact sur la fibre \(B_1\) ; le parcours nominal et le morphisme déclaré expriment le caractère scalaire sans remplacer le spectre de la fibre.
\medskip
\Needspace{10\baselineskip}
\noindent\textbf{secteur leptonique chargé, profondeur G2.} \(\mathcal B_{\ell,\mathrm{G2}}\).\par
\noindent\textbf{Domaine de fibre.} Rang 8; cellules 21,23,25,39,43,57,59,61; sans orbite chromatique; 8 familles de parcours.
\noindent\textbf{Spectre complet de \(K_D\).}
\begin{itemize}
\item \texttt{\detokenize{(0,24,0)x1}}
\item \texttt{\detokenize{(40,78,27)x2}}
\item \texttt{\detokenize{(80,54,7)x1}}
\item \texttt{\detokenize{(80,66,2)x1}}
\item \texttt{\detokenize{(80,66,3)x1}}
\item \texttt{\detokenize{(100,66,0)x1}}
\item \texttt{\detokenize{(100,66,2)x1}}
\end{itemize}
\noindent\textbf{Spectre complet de \(K_\Omega\).}
\begin{itemize}
\item \texttt{\detokenize{(80,54,6)x8}}
\end{itemize}
\noindent\textbf{Conclusion dans la fibre.} Opérateur diagonal exact sur la fibre \(B_1\) ; le parcours nominal et le morphisme déclaré expriment le caractère scalaire sans remplacer le spectre de la fibre.
\medskip
\Needspace{22\baselineskip}
\noindent\textbf{secteur leptonique chargé, profondeur G4.} \(\mathcal B_{\ell,\mathrm{G4}}\).\par
\noindent\textbf{Domaine de fibre.} Rang 16; cellules 1,3,5,7,9,19,27,37,45,55,63,73,75,77,79,81; sans orbite chromatique; 16 familles de parcours.
\noindent\textbf{Spectre complet de \(K_D\).}
\begin{itemize}
\item \texttt{\detokenize{(0,24,0)x1}}
\item \texttt{\detokenize{(40,24,2)x2}}
\item \texttt{\detokenize{(40,54,6)x2}}
\item \texttt{\detokenize{(40,84,30)x2}}
\item \texttt{\detokenize{(60,78,25)x3}}
\item \texttt{\detokenize{(80,66,2)x2}}
\item \texttt{\detokenize{(80,66,3)x3}}
\item \texttt{\detokenize{(80,66,4)x1}}
\end{itemize}
\noindent\textbf{Spectre complet de \(K_\Omega\).}
\begin{itemize}
\item \texttt{\detokenize{(24,54,2)x1}}
\item \texttt{\detokenize{(56,54,5)x4}}
\item \texttt{\detokenize{(80,54,6)x11}}
\end{itemize}
\noindent\textbf{Conclusion dans la fibre.} Opérateur diagonal exact sur la fibre \(B_1\) ; le parcours nominal et le morphisme déclaré expriment le caractère scalaire sans remplacer le spectre de la fibre.
\medskip
\Needspace{10\baselineskip}
\noindent\textbf{secteur neutrinique, profondeur G1.} \(\mathcal B_{\nu,\mathrm{G1}}\).\par
\noindent\textbf{Domaine de fibre.} Rang 2; cellules 40,42; sans orbite chromatique; 2 familles de parcours.
\noindent\textbf{Spectre complet de \(K_D\).}
\begin{itemize}
\item \texttt{\detokenize{(40,24,2)x1}}
\item \texttt{\detokenize{(40,78,27)x1}}
\end{itemize}
\noindent\textbf{Spectre complet de \(K_\Omega\).}
\begin{itemize}
\item \texttt{\detokenize{(40,54,4)x1}}
\item \texttt{\detokenize{(80,54,6)x1}}
\end{itemize}
\noindent\textbf{Conclusion dans la fibre.} Opérateur diagonal exact sur la fibre \(B_1\) ; le parcours nominal et le morphisme déclaré expriment le caractère scalaire sans remplacer le spectre de la fibre.
\medskip
\Needspace{10\baselineskip}
\noindent\textbf{secteur neutrinique, profondeur G2.} \(\mathcal B_{\nu,\mathrm{G2}}\).\par
\noindent\textbf{Domaine de fibre.} Rang 4; cellules 22,24,58,60; sans orbite chromatique; 4 familles de parcours.
\noindent\textbf{Spectre complet de \(K_D\).}
\begin{itemize}
\item \texttt{\detokenize{(0,24,0)x1}}
\item \texttt{\detokenize{(40,84,30)x1}}
\item \texttt{\detokenize{(60,78,25)x1}}
\item \texttt{\detokenize{(80,66,3)x1}}
\end{itemize}
\noindent\textbf{Spectre complet de \(K_\Omega\).}
\begin{itemize}
\item \texttt{\detokenize{(4,54,2)x1}}
\item \texttt{\detokenize{(80,54,6)x3}}
\end{itemize}
\noindent\textbf{Conclusion dans la fibre.} Opérateur diagonal exact sur la fibre \(B_1\) ; le parcours nominal et le morphisme déclaré expriment le caractère scalaire sans remplacer le spectre de la fibre.
\medskip
\Needspace{18\baselineskip}
\noindent\textbf{secteur neutrinique, profondeur G3.} \(\mathcal B_{\nu,\mathrm{G3}}\).\par
\noindent\textbf{Domaine de fibre.} Rang 6; cellules 20,26,38,44,56,62; sans orbite chromatique; 6 familles de parcours.
\noindent\textbf{Spectre complet de \(K_D\).}
\begin{itemize}
\item \texttt{\detokenize{(40,24,2)x2}}
\item \texttt{\detokenize{(40,78,27)x1}}
\item \texttt{\detokenize{(60,84,28)x1}}
\item \texttt{\detokenize{(80,54,6)x1}}
\item \texttt{\detokenize{(80,66,3)x1}}
\end{itemize}
\noindent\textbf{Spectre complet de \(K_\Omega\).}
\begin{itemize}
\item \texttt{\detokenize{(36,54,4)x1}}
\item \texttt{\detokenize{(56,54,5)x1}}
\item \texttt{\detokenize{(80,54,6)x4}}
\end{itemize}
\noindent\textbf{Conclusion dans la fibre.} Opérateur diagonal exact sur la fibre \(B_1\) ; le parcours nominal et le morphisme déclaré expriment le caractère scalaire sans remplacer le spectre de la fibre.
\medskip
\Needspace{10\baselineskip}
\noindent\textbf{secteur neutrinique, profondeur G4.} \(\mathcal B_{\nu,\mathrm{G4}}\).\par
\noindent\textbf{Domaine de fibre.} Rang 8; cellules 2,4,6,8,74,76,78,80; sans orbite chromatique; 8 familles de parcours.
\noindent\textbf{Spectre complet de \(K_D\).}
\begin{itemize}
\item \texttt{\detokenize{(0,24,0)x1}}
\item \texttt{\detokenize{(40,24,2)x1}}
\item \texttt{\detokenize{(40,78,27)x2}}
\item \texttt{\detokenize{(40,84,30)x1}}
\item \texttt{\detokenize{(80,54,7)x1}}
\item \texttt{\detokenize{(80,66,2)x1}}
\item \texttt{\detokenize{(100,66,0)x1}}
\end{itemize}
\noindent\textbf{Spectre complet de \(K_\Omega\).}
\begin{itemize}
\item \texttt{\detokenize{(36,54,4)x1}}
\item \texttt{\detokenize{(56,54,5)x1}}
\item \texttt{\detokenize{(80,54,6)x6}}
\end{itemize}
\noindent\textbf{Conclusion dans la fibre.} Opérateur diagonal exact sur la fibre \(B_1\) ; le parcours nominal et le morphisme déclaré expriment le caractère scalaire sans remplacer le spectre de la fibre.
\medskip
\Needspace{10\baselineskip}
\noindent\textbf{branche de quarks de type bas, profondeur G1.} \(\mathcal B_{d,\mathrm{G1}}\).\par
\noindent\textbf{Domaine de fibre.} Rang 2; cellules 32,50; B | G; 2 familles de parcours.
\noindent\textbf{Spectre complet de \(K_D\).}
\begin{itemize}
\item \texttt{\detokenize{(40,24,2)x1}}
\item \texttt{\detokenize{(40,78,27)x1}}
\end{itemize}
\noindent\textbf{Spectre complet de \(K_\Omega\).}
\begin{itemize}
\item \texttt{\detokenize{(40,54,4)x1}}
\item \texttt{\detokenize{(80,54,6)x1}}
\end{itemize}
\noindent\textbf{Conclusion dans la fibre.} Opérateur diagonal exact sur la fibre \(B_1\) ; le parcours nominal et le morphisme déclaré expriment le caractère scalaire sans remplacer le spectre de la fibre.
\medskip
\Needspace{10\baselineskip}
\noindent\textbf{branche de quarks de type bas, profondeur G2.} \(\mathcal B_{d,\mathrm{G2}}\).\par
\noindent\textbf{Domaine de fibre.} Rang 4; cellules 30,34,48,52; B | G | R; 4 familles de parcours.
\noindent\textbf{Spectre complet de \(K_D\).}
\begin{itemize}
\item \texttt{\detokenize{(0,24,0)x1}}
\item \texttt{\detokenize{(40,84,30)x1}}
\item \texttt{\detokenize{(60,78,25)x1}}
\item \texttt{\detokenize{(80,66,2)x1}}
\end{itemize}
\noindent\textbf{Spectre complet de \(K_\Omega\).}
\begin{itemize}
\item \texttt{\detokenize{(4,54,2)x1}}
\item \texttt{\detokenize{(80,54,6)x3}}
\end{itemize}
\noindent\textbf{Conclusion dans la fibre.} Opérateur diagonal exact sur la fibre \(B_1\) ; le parcours nominal et le morphisme déclaré expriment le caractère scalaire sans remplacer le spectre de la fibre.
\medskip
\Needspace{10\baselineskip}
\noindent\textbf{branche de quarks de type bas, profondeur G3.} \(\mathcal B_{d,\mathrm{G3}}\).\par
\noindent\textbf{Domaine de fibre.} Rang 6; cellules 12,14,16,66,68,70; B | G | R; 6 familles de parcours.
\noindent\textbf{Spectre complet de \(K_D\).}
\begin{itemize}
\item \texttt{\detokenize{(40,24,2)x2}}
\item \texttt{\detokenize{(40,78,27)x1}}
\item \texttt{\detokenize{(60,84,28)x1}}
\item \texttt{\detokenize{(80,54,6)x1}}
\item \texttt{\detokenize{(80,66,4)x1}}
\end{itemize}
\noindent\textbf{Spectre complet de \(K_\Omega\).}
\begin{itemize}
\item \texttt{\detokenize{(36,54,4)x1}}
\item \texttt{\detokenize{(80,54,6)x5}}
\end{itemize}
\noindent\textbf{Conclusion dans la fibre.} Opérateur diagonal exact sur la fibre \(B_1\) ; le parcours nominal et le morphisme déclaré expriment le caractère scalaire sans remplacer le spectre de la fibre.
\medskip
\Needspace{10\baselineskip}
\noindent\textbf{branche de quarks de type bas, profondeur G4.} \(\mathcal B_{d,\mathrm{G4}}\).\par
\noindent\textbf{Domaine de fibre.} Rang 8; cellules 10,18,28,36,46,54,64,72; B | G | R; 8 familles de parcours.
\noindent\textbf{Spectre complet de \(K_D\).}
\begin{itemize}
\item \texttt{\detokenize{(0,24,0)x1}}
\item \texttt{\detokenize{(40,24,2)x1}}
\item \texttt{\detokenize{(40,78,27)x2}}
\item \texttt{\detokenize{(40,84,30)x1}}
\item \texttt{\detokenize{(80,54,7)x1}}
\item \texttt{\detokenize{(80,66,3)x1}}
\item \texttt{\detokenize{(100,66,2)x1}}
\end{itemize}
\noindent\textbf{Spectre complet de \(K_\Omega\).}
\begin{itemize}
\item \texttt{\detokenize{(36,54,4)x1}}
\item \texttt{\detokenize{(80,54,6)x7}}
\end{itemize}
\noindent\textbf{Conclusion dans la fibre.} Opérateur diagonal exact sur la fibre \(B_1\) ; le parcours nominal et le morphisme déclaré expriment le caractère scalaire sans remplacer le spectre de la fibre.
\medskip
\Needspace{10\baselineskip}
\noindent\textbf{branche de quarks de type haut, profondeur G1.} \(\mathcal B_{u,\mathrm{G1}}\).\par
\noindent\textbf{Domaine de fibre.} Rang 4; cellules 31,33,49,51; B | G | R; 3 familles de parcours.
\noindent\textbf{Spectre complet de \(K_D\).}
\begin{itemize}
\item \texttt{\detokenize{(40,54,6)x1}}
\item \texttt{\detokenize{(60,78,25)x1}}
\item \texttt{\detokenize{(80,66,2)x1}}
\item \texttt{\detokenize{(80,66,3)x1}}
\end{itemize}
\noindent\textbf{Spectre complet de \(K_\Omega\).}
\begin{itemize}
\item \texttt{\detokenize{(80,54,6)x4}}
\end{itemize}
\noindent\textbf{Conclusion dans la fibre.} Opérateur diagonal exact sur la fibre \(B_1\) ; le parcours nominal et le morphisme déclaré expriment le caractère scalaire sans remplacer le spectre de la fibre.
\medskip
\Needspace{29\baselineskip}
\noindent\textbf{branche de quarks de type haut, profondeur G3.} \(\mathcal B_{u,\mathrm{G3}}\).\par
\noindent\textbf{Domaine de fibre.} Rang 12; cellules 11,13,15,17,29,35,47,53,65,67,69,71; B | G | R; 12 familles de parcours.
\noindent\textbf{Spectre complet de \(K_D\).}
\begin{itemize}
\item \texttt{\detokenize{(40,24,2)x2}}
\item \texttt{\detokenize{(40,78,27)x2}}
\item \texttt{\detokenize{(60,84,28)x1}}
\item \texttt{\detokenize{(80,54,6)x3}}
\item \texttt{\detokenize{(80,66,3)x1}}
\item \texttt{\detokenize{(80,66,4)x1}}
\item \texttt{\detokenize{(100,66,0)x1}}
\item \texttt{\detokenize{(100,66,2)x1}}
\end{itemize}
\noindent\textbf{Spectre complet de \(K_\Omega\).}
\begin{itemize}
\item \texttt{\detokenize{(56,54,5)x3}}
\item \texttt{\detokenize{(80,54,6)x9}}
\end{itemize}
\noindent\textbf{Conclusion dans la fibre.} Opérateur diagonal exact sur la fibre \(B_1\) ; le parcours nominal et le morphisme déclaré expriment le caractère scalaire sans remplacer le spectre de la fibre.
\medskip
\endgroup


\Needspace{15\baselineskip}
\subsubsection{Morphismes de réalisation physique}

\begingroup
\setlength{\tabcolsep}{4.5pt}
\renewcommand{\arraystretch}{1.28}
\begin{longtable}{@{}>{\raggedright\arraybackslash}p{.14\textwidth}>{\raggedright\arraybackslash}p{.33\textwidth}>{\raggedright\arraybackslash}p{.25\textwidth}>{\raggedright\arraybackslash}p{.22\textwidth}@{}}
\caption{Réalisation physique des leptons chargés.}\label{tab:masas-realizacion-leptones}\\
\toprule
\textbf{Classe} & \textbf{Objeto HMT} & \textbf{Application de réalisation} & \textbf{Conclusion et conditions} \\
\midrule
\endfirsthead
\toprule
\textbf{Classe} & \textbf{Objeto HMT} & \textbf{Application de réalisation} & \textbf{Conclusion et conditions} \\
\midrule
\endhead
\midrule
\multicolumn{4}{r}{\footnotesize Suite à la page suivante.}\\
\endfoot
\bottomrule
\endlastfoot
$e^-/e^+$ & Cellule centrale $(5,5)$, parcours $\Gamma_e=(5,5,++)$ et $K_D(\Gamma_e)=K_\Omega(\Gamma_e)=(80,54,6)$. & Réduction de rang un fermée. & Clôture bêta de rang un ; sélecteur indépendant des grandeurs métrologiques utilisées comme cibles. \\
$\mu^-/\mu^+$ & famille leptonique chargée ; niveaux G2 ; 8 cellules ; $K_D$ : 7 profils ; $K_\Omega$ : 1 profil & Génération et signature fine $\to$ parcours persistant $\to\chi_{\rm mass}\to\mathcal L_M$. & Parcours nominal et caractère scalaire fermés ; la paire d’opérateurs conserve le raffinement complet de la fibre. \\
$\tau^-/\tau^+$ & famille leptonique chargée ; niveaux G4 ; 16 cellules ; $K_D$ : 8 profils ; $K_\Omega$ : 3 profils & Génération et signature fine $\to$ parcours persistant $\to\chi_{\rm mass}\to\mathcal L_M$. & Parcours nominal et caractère scalaire fermés ; la paire d’opérateurs conserve le raffinement complet de la fibre. \\
\end{longtable}
\endgroup

\begingroup
\setlength{\tabcolsep}{4.5pt}
\renewcommand{\arraystretch}{1.28}
\begin{longtable}{@{}>{\raggedright\arraybackslash}p{.14\textwidth}>{\raggedright\arraybackslash}p{.33\textwidth}>{\raggedright\arraybackslash}p{.25\textwidth}>{\raggedright\arraybackslash}p{.22\textwidth}@{}}
\caption{Réalisation physique des neutrinos et transport de bases.}\label{tab:masas-realizacion-neutrinos}\\
\toprule
\textbf{Classe} & \textbf{Objeto HMT} & \textbf{Application de réalisation} & \textbf{Conclusion et conditions} \\
\midrule
\endfirsthead
\toprule
\textbf{Classe} & \textbf{Objeto HMT} & \textbf{Application de réalisation} & \textbf{Conclusion et conditions} \\
\midrule
\endhead
\midrule
\multicolumn{4}{r}{\footnotesize Suite à la page suivante.}\\
\endfoot
\bottomrule
\endlastfoot
$\nu_e/\bar\nu_e$ & famille neutre ; niveaux G1, G2, G3, G4 ; 20 cellules ; $K_D$ : 11 profils ; $K_\Omega$ : 5 profils & Enregistrement neutre de rang plein $\to$ vecteurs propres $\to$ base PMNS. & La structure extérieure/CAR fixe le caractère fermionique ; PMNS transporte les bases et ne crée pas de valeurs propres de masse. \\
$\nu_\mu/\bar\nu_\mu$ & famille neutre ; niveaux G1, G2, G3, G4 ; 20 cellules ; $K_D$ : 11 profils ; $K_\Omega$ : 5 profils & Enregistrement neutre de rang plein $\to$ vecteurs propres $\to$ base PMNS. & La structure extérieure/CAR fixe le caractère fermionique ; PMNS transporte les bases et ne crée pas de valeurs propres de masse. \\
$\nu_\tau/\bar\nu_\tau$ & famille neutre ; niveaux G1, G2, G3, G4 ; 20 cellules ; $K_D$ : 11 profils ; $K_\Omega$ : 5 profils & Enregistrement neutre de rang plein $\to$ vecteurs propres $\to$ base PMNS. & La structure extérieure/CAR fixe le caractère fermionique ; PMNS transporte les bases et ne crée pas de valeurs propres de masse. \\
\end{longtable}
\endgroup

\begingroup
\setlength{\tabcolsep}{4.5pt}
\renewcommand{\arraystretch}{1.28}
\begin{longtable}{@{}>{\raggedright\arraybackslash}p{.14\textwidth}>{\raggedright\arraybackslash}p{.33\textwidth}>{\raggedright\arraybackslash}p{.25\textwidth}>{\raggedright\arraybackslash}p{.22\textwidth}@{}}
\caption{Réalisation physique des quarks : génération, saveur et couleur.}\label{tab:masas-realizacion-quarks}\\
\toprule
\textbf{Classe} & \textbf{Objeto HMT} & \textbf{Application de réalisation} & \textbf{Conclusion et conditions} \\
\midrule
\endfirsthead
\toprule
\textbf{Classe} & \textbf{Objeto HMT} & \textbf{Application de réalisation} & \textbf{Conclusion et conditions} \\
\midrule
\endhead
\midrule
\multicolumn{4}{r}{\footnotesize Suite à la page suivante.}\\
\endfoot
\bottomrule
\endlastfoot
$u/\bar u$ & branche de quarks \emph{up} ; niveaux G1 et G3 ; couleur $\mathbb Z/3$ ; 16 cellules ; $K_D$ : 11 profils ; $K_\Omega$ : 2 profils & Saveur--couleur--génération $\to$ parcours persistant $\to\chi_{\rm mass}\to\mathcal L_M$. & La saveur, la couleur et la génération déterminent la classe de parcours ; caractère scalaire et paire d’opérateurs fermés. \\
$d/\bar d$ & branche de quarks \emph{down} ; niveaux G1, G2, G3, G4 ; couleur $\mathbb Z/3$ ; 20 cellules ; $K_D$ : 12 profils ; $K_\Omega$ : 4 profils & Saveur--couleur--génération $\to$ parcours persistant $\to\chi_{\rm mass}\to\mathcal L_M$. & La saveur, la couleur et la génération déterminent la classe de parcours ; caractère scalaire et paire d’opérateurs fermés. \\
$c/\bar c$ & branche de quarks \emph{up} ; niveaux G1 et G3 ; couleur $\mathbb Z/3$ ; 16 cellules ; $K_D$ : 11 profils ; $K_\Omega$ : 2 profils & Saveur--couleur--génération $\to$ parcours persistant $\to\chi_{\rm mass}\to\mathcal L_M$. & La saveur, la couleur et la génération déterminent la classe de parcours ; caractère scalaire et paire d’opérateurs fermés. \\
$s/\bar s$ & branche de quarks \emph{down} ; niveaux G1, G2, G3, G4 ; couleur $\mathbb Z/3$ ; 20 cellules ; $K_D$ : 12 profils ; $K_\Omega$ : 4 profils & Saveur--couleur--génération $\to$ parcours persistant $\to\chi_{\rm mass}\to\mathcal L_M$. & La saveur, la couleur et la génération déterminent la classe de parcours ; caractère scalaire et paire d’opérateurs fermés. \\
$t/\bar t$ & branche de quarks \emph{up} ; niveaux G1 et G3 ; couleur $\mathbb Z/3$ ; 16 cellules ; $K_D$ : 11 profils ; $K_\Omega$ : 2 profils & Saveur--couleur--génération $\to$ parcours persistant $\to\chi_{\rm mass}\to\mathcal L_M$. & La saveur, la couleur et la génération déterminent la classe de parcours ; caractère scalaire et paire d’opérateurs fermés. \\
$b/\bar b$ & branche de quarks \emph{down} ; niveaux G1, G2, G3, G4 ; couleur $\mathbb Z/3$ ; 20 cellules ; $K_D$ : 12 profils ; $K_\Omega$ : 4 profils & Saveur--couleur--génération $\to$ parcours persistant $\to\chi_{\rm mass}\to\mathcal L_M$. & La saveur, la couleur et la génération déterminent la classe de parcours ; caractère scalaire et paire d’opérateurs fermés. \\
\end{longtable}
\endgroup

\Needspace{12\baselineskip}
\paragraph{Localisation structurelle des \(6\) saveurs de quarks}
\begingroup
\setlength{\parskip}{.3\baselineskip}
La position d’une saveur de quark est spécifiée par quatre données simultanées : branche, génération, support cellulaire et famille de parcours orientés. Les trois générations d’une même branche partagent le support de l’atlas ; le classificateur fin de génération conserve la distinction physique qui doit encore être réalisée comme vecteur propre et section de masse.
\begingroup
\setlength{\tabcolsep}{4.5pt}
\renewcommand{\arraystretch}{1.28}
\begin{longtable}{@{}>{\raggedright\arraybackslash}p{.12\textwidth}>{\raggedright\arraybackslash}p{.17\textwidth}>{\raggedright\arraybackslash}p{.19\textwidth}>{\raggedright\arraybackslash}p{.46\textwidth}@{}}
\caption{Classification générationnelle des six saveurs de quarks.}\label{tab:masas-posicion-quarks-seis-sabores}\\
\toprule
\textbf{Saveur} & \textbf{Génération} & \textbf{Branche} & \textbf{Classificateur de branche et de génération} \\
\midrule
\endfirsthead
\toprule
\textbf{Saveur} & \textbf{Génération} & \textbf{Branche} & \textbf{Classificateur de branche et de génération} \\
\midrule
\endhead
\midrule
\multicolumn{4}{r}{\footnotesize Suite à la page suivante.}\\
\endfoot
\bottomrule
\endlastfoot
$u/\bar u$ & première & supérieur & $(\operatorname{up},1,\chi_{\rm fina})$ ; paire nominale de parcours $(21,30)$ sur l’axe $\mathrm{N\!-\!S}$ \\
$d/\bar d$ & première & inférieur & $(\operatorname{down},1,\chi_{\rm fina})$ ; paire nominale de parcours $(20,32)$ sur l’axe $\mathrm{E\!-\!O}$ \\
$c/\bar c$ & deuxième & supérieur & $(\operatorname{up},2,\chi_{\rm fina})$ ; paire nominale de parcours $(21,29)$ sur l’axe $\mathrm{E\!-\!O}$ \\
$s/\bar s$ & deuxième & inférieur & $(\operatorname{down},2,\chi_{\rm fina})$ ; paire nominale de parcours $(17,27)$ sur l’axe $\mathrm{N\!-\!S}$ \\
$t/\bar t$ & troisième & supérieur & $(\operatorname{up},3,\chi_{\rm fina})$ ; paire nominale de parcours $(21,26)$ sur l’axe $\mathrm{N\!-\!S}$ \\
$b/\bar b$ & troisième & inférieur & $(\operatorname{down},3,\chi_{\rm fina})$ ; paire nominale de parcours $(21,30)$ sur l’axe $\mathrm{E\!-\!O}$ \\
\end{longtable}
\endgroup

\Needspace{18\baselineskip}
\noindent\textbf{Support structurel de la branche supérieure.}\par
\noindent\textit{Multisections.} \(\mathcal B_{u,\mathrm{G1}}\), \(\mathcal B_{u,\mathrm{G3}}\).\par
\noindent\textit{Profondeurs et cardinaux.} $\mathrm{G1}$, $\mathrm{G3}$ ; 16 cellules et 14 familles de parcours distinctes.\par
\noindent\textit{Orbite chromatique et orientations.} $\mathcal C=\{B,G,R\}$ et $\mathcal O=\{(-1,-1),(-1,1),(1,-1),(1,1)\}$.\par
\noindent\textit{Cellules et familles orientées $B_1$.} Chaque chaîne $k\leftrightarrow(r,c)\leftrightarrow\Gamma_{r,c}^{(h,v)}$ conserve l’indice cellulaire, sa coordonnée dans la carte $9\times9$ et les deux signes d’orientation de l’enregistrement de parcours enrichi.\par
\noindent Fila $r=2$: $11\leftrightarrow(2,2)\leftrightarrow\Gamma_{2,2}^{(-,+)}$, $13\leftrightarrow(2,4)\leftrightarrow\Gamma_{2,4}^{(+,-)}$, $15\leftrightarrow(2,6)\leftrightarrow\Gamma_{2,6}^{(-,+)}$, $17\leftrightarrow(2,8)\leftrightarrow\Gamma_{2,8}^{(-,-)}$.\par
\noindent Fila $r=4$: $29\leftrightarrow(4,2)\leftrightarrow\Gamma_{4,2}^{(+,-)}$, $31\leftrightarrow(4,4)\leftrightarrow\Gamma_{4,4}^{(-,+)}$, $33\leftrightarrow(4,6)\leftrightarrow\Gamma_{4,6}^{(+,+)}$, $35\leftrightarrow(4,8)\leftrightarrow\Gamma_{4,8}^{(-,-)}$.\par
\noindent Fila $r=6$: $47\leftrightarrow(6,2)\leftrightarrow\Gamma_{6,2}^{(-,+)}$, $49\leftrightarrow(6,4)\leftrightarrow\Gamma_{6,4}^{(-,-)}$, $51\leftrightarrow(6,6)\leftrightarrow\Gamma_{6,6}^{(+,+)}$, $53\leftrightarrow(6,8)\leftrightarrow\Gamma_{6,8}^{(+,+)}$.\par
\noindent Fila $r=8$: $65\leftrightarrow(8,2)\leftrightarrow\Gamma_{8,2}^{(+,+)}$, $67\leftrightarrow(8,4)\leftrightarrow\Gamma_{8,4}^{(+,+)}$, $69\leftrightarrow(8,6)\leftrightarrow\Gamma_{8,6}^{(-,-)}$, $71\leftrightarrow(8,8)\leftrightarrow\Gamma_{8,8}^{(+,+)}$.\par
\noindent\textit{Morphisme de réalisation fermé.} Pour $g\in\{1,2,3\}$, la branche, la génération, l’orbite de couleur et la signature fine déterminent
\[\chi_{\rm fina}(\operatorname{up},g,\mathcal C)\longrightarrow\Gamma_{\rm enr}\longrightarrow\bigl(\operatorname{quot}_{\mathcal C},S(B_D,B_\Omega)\bigr).\]
Cette flèche sélectionne le vecteur propre physique, puis exprime la coordonnée scalaire dans $\mathcal L_M$ sans utiliser la valeur de comparaison.\par
\medskip
\Needspace{18\baselineskip}
\noindent\textbf{Support structurel de la branche inférieure.}\par
\noindent\textit{Multisections.} \(\mathcal B_{d,\mathrm{G1}}\), \(\mathcal B_{d,\mathrm{G2}}\), \(\mathcal B_{d,\mathrm{G3}}\), \(\mathcal B_{d,\mathrm{G4}}\).\par
\noindent\textit{Profondeurs et cardinaux.} $\mathrm{G1}$, $\mathrm{G2}$, $\mathrm{G3}$, $\mathrm{G4}$ ; 20 cellules et 18 familles de parcours distinctes.\par
\noindent\textit{Orbite chromatique et orientations.} $\mathcal C=\{B,G,R\}$ et $\mathcal O=\{(-1,-1),(-1,1),(1,-1),(1,1)\}$.\par
\noindent\textit{Cellules et familles orientées $B_1$.} Chaque chaîne $k\leftrightarrow(r,c)\leftrightarrow\Gamma_{r,c}^{(h,v)}$ conserve l’indice cellulaire, sa coordonnée dans la carte $9\times9$ et les deux signes d’orientation de l’enregistrement de parcours enrichi.\par
\noindent Fila $r=2$: $10\leftrightarrow(2,1)\leftrightarrow\Gamma_{2,1}^{(+,+)}$, $12\leftrightarrow(2,3)\leftrightarrow\Gamma_{2,3}^{(-,-)}$, $14\leftrightarrow(2,5)\leftrightarrow\Gamma_{2,5}^{(+,+)}$, $16\leftrightarrow(2,7)\leftrightarrow\Gamma_{2,7}^{(-,+)}$, $18\leftrightarrow(2,9)\leftrightarrow\Gamma_{2,9}^{(-,+)}$.\par
\noindent Fila $r=4$: $28\leftrightarrow(4,1)\leftrightarrow\Gamma_{4,1}^{(-,-)}$, $30\leftrightarrow(4,3)\leftrightarrow\Gamma_{4,3}^{(+,+)}$, $32\leftrightarrow(4,5)\leftrightarrow\Gamma_{4,5}^{(+,-)}$, $34\leftrightarrow(4,7)\leftrightarrow\Gamma_{4,7}^{(-,+)}$, $36\leftrightarrow(4,9)\leftrightarrow\Gamma_{4,9}^{(-,+)}$.\par
\noindent Fila $r=6$: $46\leftrightarrow(6,1)\leftrightarrow\Gamma_{6,1}^{(-,-)}$, $48\leftrightarrow(6,3)\leftrightarrow\Gamma_{6,3}^{(-,+)}$, $50\leftrightarrow(6,5)\leftrightarrow\Gamma_{6,5}^{(-,+)}$, $52\leftrightarrow(6,7)\leftrightarrow\Gamma_{6,7}^{(-,+)}$, $54\leftrightarrow(6,9)\leftrightarrow\Gamma_{6,9}^{(+,-)}$.\par
\noindent Fila $r=8$: $64\leftrightarrow(8,1)\leftrightarrow\Gamma_{8,1}^{(+,-)}$, $66\leftrightarrow(8,3)\leftrightarrow\Gamma_{8,3}^{(-,+)}$, $68\leftrightarrow(8,5)\leftrightarrow\Gamma_{8,5}^{(-,+)}$, $70\leftrightarrow(8,7)\leftrightarrow\Gamma_{8,7}^{(+,-)}$, $72\leftrightarrow(8,9)\leftrightarrow\Gamma_{8,9}^{(-,+)}$.\par
\noindent\textit{Morphisme de réalisation fermé.} Pour $g\in\{1,2,3\}$, la branche, la génération, l’orbite de couleur et la signature fine déterminent
\[\chi_{\rm fina}(\operatorname{down},g,\mathcal C)\longrightarrow\Gamma_{\rm enr}\longrightarrow\bigl(\operatorname{quot}_{\mathcal C},S(B_D,B_\Omega)\bigr).\]
Cette flèche sélectionne le vecteur propre physique, puis exprime la coordonnée scalaire dans $\mathcal L_M$ sans utiliser la valeur de comparaison.\par
\medskip
\endgroup

\begingroup
\setlength{\tabcolsep}{4.5pt}
\renewcommand{\arraystretch}{1.28}
\begin{longtable}{@{}>{\raggedright\arraybackslash}p{.14\textwidth}>{\raggedright\arraybackslash}p{.33\textwidth}>{\raggedright\arraybackslash}p{.25\textwidth}>{\raggedright\arraybackslash}p{.22\textwidth}@{}}
\caption{Réalisation physique des secteurs de jauge et scalaire.}\label{tab:masas-realizacion-gauge-escalar}\\
\toprule
\textbf{Classe} & \textbf{Objeto HMT} & \textbf{Application de réalisation} & \textbf{Conclusion et conditions} \\
\midrule
\endfirsthead
\toprule
\textbf{Classe} & \textbf{Objeto HMT} & \textbf{Application de réalisation} & \textbf{Conclusion et conditions} \\
\midrule
\endhead
\midrule
\multicolumn{4}{r}{\footnotesize Suite à la page suivante.}\\
\endfoot
\bottomrule
\endlastfoot
photon $\gamma$ & Carte nulle sectorielle et représentation de jauge correspondante. & Inclusion exacte de la section nulle dans la droite de masse. & Nullité exacte dans sa carte sectorielle. \\
gluones $g$ & Carte nulle sectorielle et représentation de jauge correspondante. & Inclusion exacte de la section nulle dans la droite de masse. & Nullité exacte dans sa carte sectorielle. \\
$W^\pm$ & parcours CT--108 W : 15/37, axe E-O, $(N,E,S,O)=(23,44,11,30)$, feuillet/résidu $(41,8)$, $(n_A,s_C)=(94,-1)$, $(W_+,W_-)=(563,565)$ ; caractère scalaire exact & Parcours CT--108 nominal $\to(n_A,s_C)\to(W_+,W_-)\to\chi_{\rm mass}\to\mathcal L_M$. & Parcours de jauge CT--108, signature et caractère scalaire déterminés avant la comparaison. \\
$Z^0$ & parcours CT--108 Z : 20/23, axe E-O, $(N,E,S,O)=(33,45,10,20)$, feuillet/résidu $(57,7)$, $(n_A,s_C)=(95,-1)$, $(W_+,W_-)=(569,571)$ ; caractère scalaire exact & Parcours CT--108 nominal $\to(n_A,s_C)\to(W_+,W_-)\to\chi_{\rm mass}\to\mathcal L_M$. & Parcours de jauge CT--108, signature et caractère scalaire déterminés avant la comparaison. \\
Higgs $H$ & parcours CT--108 H : 24/29, axe N-S, $(N,E,S,O)=(40,34,22,12)$, feuillet/résidu $(47,10)$, $(n_A,s_C)=(97,9)$, $(W_+,W_-)=(591,573)$ ; caractère scalaire exact & Parcours CT--108 nominal $\to(n_A,s_C)\to(W_+,W_-)\to\chi_{\rm mass}\to\mathcal L_M$. & Parcours scalaire CT--108, signature et caractère déterminés avant la comparaison. \\
\end{longtable}
\endgroup

\begingroup
\setlength{\tabcolsep}{4.5pt}
\renewcommand{\arraystretch}{1.28}
\begin{longtable}{@{}>{\raggedright\arraybackslash}p{.14\textwidth}>{\raggedright\arraybackslash}p{.33\textwidth}>{\raggedright\arraybackslash}p{.25\textwidth}>{\raggedright\arraybackslash}p{.22\textwidth}@{}}
\caption{Réalisation des composés, des secteurs topologiques et des états virtuels.}\label{tab:masas-realizacion-no-escalar}\\
\toprule
\textbf{Classe} & \textbf{Objeto HMT} & \textbf{Application de réalisation} & \textbf{Conclusion et conditions} \\
\midrule
\endfirsthead
\toprule
\textbf{Classe} & \textbf{Objeto HMT} & \textbf{Application de réalisation} & \textbf{Conclusion et conditions} \\
\midrule
\endhead
\midrule
\multicolumn{4}{r}{\footnotesize Suite à la page suivante.}\\
\endfoot
\bottomrule
\endlastfoot
mésons & Enregistrement double, frontière de liaison et signature composée. & Enregistrement composé et liaison $\to$ carte scalaire. & Composition des enregistrements ; la liaison précède la carte scalaire. \\
baryons & Enregistrement triple, retenue et signature composée. & Enregistrement composé et liaison $\to$ carte scalaire. & Composition des enregistrements ; la liaison précède la carte scalaire. \\
Fibonacci $(1,\tau)$ & Anneau de fusion de Fibonacci : $\tau\otimes\tau=1+\tau$. & Fusion et tressage $\to$ observable physique du secteur. & Codomaine exact de fusion et de tressage ; la comparaison utilise ses invariants propres. \\
Ising $(1,\psi,\sigma)$ ; réalisation de Majorana séparée & Anneau de fusion d’Ising : $(1,\psi,\sigma)$. & Fusion et tressage $\to$ observable physique du secteur. & Codomaine exact de fusion et de tressage ; la comparaison utilise ses invariants propres. \\
sectores $\mathbb Z/6$ & Caractère de tressage sur $\mathbb Z/6$. & Fusion et tressage $\to$ observable physique du secteur. & Codomaine exact de fusion et de tressage ; la comparaison utilise ses invariants propres. \\
secciones virtuales & Système inverse d’extensions et horizon de persistance. & Prolongement enregistré $\to$ horizon physique de persistance. & Objet exact de prolongement et d’obstruction ; sa sortie est l’horizon de persistance. \\
\end{longtable}
\endgroup


\Needspace{15\baselineskip}
\subsubsection{Comparabilité physique dans les codomaines natifs}

La comparaison externe se définit sur le codomaine natif de chaque classe. Dans les
classes à section scalaire établie, on compare des grandeurs homogènes, avec unité,
schéma et échelle déclarés. Dans les fibres opératorielles, on compare le spectre
interne et l'on conserve explicitement le morphisme de réalisation établi qui
sélectionne le parcours sans utiliser la valeur externe. Les enregistrements mésoniques et baryoniques n'admettent une comparaison état par
état qu'après incorporation de la frontière de liaison. Dans le secteur de Fibonacci,
l'anneau basé et \(\operatorname{FPdim}(\tau)=\varphi\) sont établis ; les matrices
\(F/R\), le pentagone, l'hexagone, le Hamiltonien et le gap relèvent de la
réalisation ultérieure. Ising et \(\mathbb Z/6\) conservent séparément leurs invariants
de fusion et de tressage, et la virtualité se compare par la persistance et la structure
analytique. Dans aucun de ces six
cas il n'existe de masse scalaire universelle de classe ; en attribuer une remplacerait l'objet
démontré par une valeur numérique sans morphisme.
